% Generated draft: requires mathematical and source review.
\documentclass[11pt]{article}
\usepackage[margin=25mm]{geometry}
\usepackage{fontspec}
\setmainfont{Latin Modern Roman}
\usepackage{amsmath,amssymb,mathtools,unicode-math}
\setmathfont{Latin Modern Math}
\usepackage{enumitem,needspace,xurl,hyperref}
\setlength{\parindent}{0pt}
\setlength{\parskip}{5pt}
\setlength{\emergencystretch}{3em}
\newcommand{\N}{\mathbb{N}}
\newcommand{\Z}{\mathbb{Z}}
\newcommand{\Q}{\mathbb{Q}}
\newcommand{\R}{\mathbb{R}}
\newcommand{\C}{\mathbb{C}}
\newcommand{\F}{\mathbb{F}}
\newcommand{\A}{\mathbb{A}}
\newcommand{\PP}{\mathbb{P}}
\newcommand{\E}{\mathbb{E}}
\newcommand{\eps}{\varepsilon}
\newcommand{\dd}{\,\mathrm{d}}
\begin{document}
\section*{UnsolvedMath ID: 20002827}
\textbf{Source problem code:} catalogue entry
\subsection{Definitions and mathematical statement}
Let $q$ be a prime power, $n\geq1$, $G=\operatorname{GL}_n(\F_q)$ and $U=U_n(q)$ its upper unitriangular subgroup. The original question asks both for relations between the supercharacter theory of $U$ and properties of $G$, and for analogous problems in which such a relation has been worked out. No particular property or analogy is prescribed.

A concrete affirmative relation reported by the supplied research aid is the following. Let $V_n(q)$ be the complex standard Diaconis--Isaacs supercharacter space of $U$, equivalently the space spanned by its standard superclass indicators. Its basis is indexed by set partitions $\pi$ of $[n]$ with nonzero $\F_q$-labels on their arcs. If $|\pi|$ is the number of blocks, define
\[
L_n(q)=\sum_{\pi\text{ a set partition of }[n]}(q-1)^{n-|\pi|}=\dim V_n(q).
\]
Let $W_n(q)$ be the space of complex class functions on $G$ supported on unipotent elements, and let $p(n)$ be the number of integer partitions of $n$. Since the unipotent conjugacy classes of $G$ are indexed by those partitions, $\dim W_n(q)=p(n)$. Ordinary induction gives
\[
I:V_n(q)\longrightarrow W_n(q),\qquad
I(f)(g)=\frac1{|U|}\sum_{\substack{x\in G\\x^{-1}gx\in U}}f(x^{-1}gx).
\]
The aid asserts that $I$ is surjective, hence has rank $p(n)$ and
\[
\dim\ker I=L_n(q)-p(n).
\]

More precisely, let $T\subset G$ be the diagonal torus, acting on $V_n(q)$ by $(t\cdot f)(u)=f(t^{-1}ut)$, and write
\[
V_n(q)_T=V_n(q)/\operatorname{span}_{\C}\{t\cdot f-f:t\in T,\ f\in V_n(q)\}.
\]
Induction factors through this coinvariant space. The reported torus orbits forget all nonzero arc labels and retain the underlying set partition, so
\[
\dim V_n(q)_T=B_n,\qquad
\overline I:V_n(q)_T\twoheadrightarrow W_n(q),\qquad
\dim\ker\overline I=B_n-p(n),
\]
where $B_n$ is the $n$th Bell number, not the Borel subgroup notation sometimes also used in this workshop.

The supplied literature explanation is that suitable nonnesting supercharacters of $U_n(q)$ induce to the generalized Gelfand--Graev basis of unipotently supported class functions of $G$, and that their multiplicities against unipotent characters are described by Kostka--Foulkes polynomials. The exact nonnesting conventions, polynomial normalization and bibliographic source are not supplied, so an explicit multiplicity identity is not asserted here. The quotient and kernel dimensions are the aid's stated partial result; they do not specify all possible relations or analogies requested in the broad workshop question.
\subsection{Short English statement}
Relate unitriangular supercharacters to representation theory of $\operatorname{GL}_n(q)$, and identify worked-out analogues. The supplied aid gives ordinary induction onto unipotently supported class functions, with rank $p(n)$ and kernel dimensions before and after diagonal-torus coinvariants.
\subsection{Sources}
\begin{itemize}
\item Local UnsolvedMath registry, permanent ID $20002827$, source code \texttt{catalogue entry}; canonical \texttt{aim-representation-theory-notes.json}, \texttt{notes[98]}.
\item AIM Problem Lists, Representation theory, workshop \emph{Supercharacters and combinatorial Hopf algebras}, Problem 1.15: \url{https://aimath.org/WWN/supercharacters/supercharacters.pdf}. The extraction has no remarks or literature field. It records the question as exact; the statement-verification field ends with ``The AIM workshop PDF says:'' and provides no further quotation.
\item The supplied research summary refers to post-workshop literature, nonnesting supercharacters, generalized Gelfand--Graev characters, Kostka--Foulkes polynomials, and the standard Diaconis--Isaacs supercharacter theory, but supplies no authoritative bibliographic citation for the induction theorem.
\item Research attempt 1: \texttt{partial\_result}, status \texttt{partially\_solved}, classification \texttt{PARTIAL-PROGRESS}, rigor label \texttt{proved}, novelty confidence low. This machine-generated aid is not peer reviewed; its statements require independent verification before citation.
\end{itemize}
\end{document}
